\section{Extension to an external magnetic field}\label{sec:magnetic}

We now add a deterministic external field. This is a separate extension of
Theorem~\ref{thm:main}. Its proof uses magnetization constraints: the
single-site functional is concave in covariance levels after prescribing
the mean spin, and a product of large constrained blocks supplies the
corresponding cavity lower bound.

For a deterministic vector $c^{(N)}=(c_1^{(N)},\ldots,c_N^{(N)})\in\R^N$,
write
\begin{equation}\label{mag:pressure}
 P_N(c^{(N)})=\frac1N\log\left[
  2^{-N}\sum_{\sigma\in\{-1,1\}^N}
  \exp\left\{\frac12\sigma^{\mathsf T}U^{\mathsf T}\Lambda_NU\sigma
                    +\sum_{r=1}^N c_r^{(N)}\sigma_r\right\}\right].
\end{equation}
For $b\in\R$, let $\ell_b^{\mathrm{can}}(h)$ denote the single-site
functional defined earlier, with initial position $b$ in place of zero;
the subtraction $h_{k-1}/2$ is unchanged. For $s\in(-1,1)$ define
\begin{equation}\label{mag:constrained-functional}
 \ell^{[s]}(h)=\inf_{b\in\R}
       \{\ell_b^{\mathrm{can}}(h)-bs\}.
\end{equation}
The bias $b$ is optimized to prescribe the mean spin $s$; the physical
external field will contribute its linear term separately in the pressure
formula. This single-site concave conjugate is used in generalized TAP theory
\cite[\S2.3, equations~(2.21)--(2.23)]{ChenPanchenkoSubag2023}.
Optimization over a conventional order parameter outside the
Parisi order parameter also appears in Gaussian vector-spin formulas
\cite[Theorem~1.1]{ChenConventional2024}.
The covariance-level concavity and the constrained-block comparison
needed for the present spectral model are proved in this section.

Given positive rational numbers $\gamma_1,\ldots,\gamma_J$ with sum one
and $\mathbf m=(m_1,\ldots,m_J)\in((-1,1)\cap\mathbb Q)^J$, put
\begin{equation}\label{mag:Sm}
 \ell^{[\mathbf m]}(h)=\sum_{j=1}^J\gamma_j\ell^{[m_j]}(h),
 \qquad
 S_{\mathbf m}(p)=\sup_{h\in\cH}
       \left\{\ell^{[\mathbf m]}(h)+\frac12\int_0^1p(s)h(s)\,ds\right\}.
\end{equation}
The same covariance path $h$ is used in every field group.

For probability laws $\nu,\widetilde\nu$ on $\R$ with finite absolute
first moments, write
\[
 W_1(\nu,\widetilde\nu)
   =\inf_{\pi}\int_{\R^2}|x-y|\,\pi(dx,dy),
\]
where the infimum is over probability laws $\pi$ with marginals $\nu$
and $\widetilde\nu$.

\begin{theorem}[Magnetic-field extension]\label{mag:theorem}
Assume the deterministic spectral hypotheses of Theorem~\ref{thm:main},
with limiting compactly supported law $\mu$. Suppose first that the field
values belong to a fixed finite alphabet $c_1,\ldots,c_J\in\R$, and that
the proportion of sites with value $c_j$ tends to a positive rational
number $\gamma_j$. Then $P_N(c^{(N)})$ converges in expectation and almost
surely to
\begin{equation}\label{mag:finite-formula}
 \mathcal F_\mu\left(\sum_{j=1}^J\gamma_j\delta_{c_j}\right)
 =\sup_{\mathbf m\in((-1,1)\cap\mathbb Q)^J}
   \left\{\sum_{j=1}^J\gamma_jc_jm_j
    +\inf_{p\in\cP}\bigl[S_{\mathbf m}(p)+G_1(p)\bigr]\right\}.
\end{equation}
More generally, suppose the empirical field laws converge in the
first-moment transport metric $W_1$ to a probability law $\nu$ on $\R$
with finite absolute first moment. Then the same modes of convergence
hold, with limit $\mathcal F_\mu(\nu)$, where
\begin{equation}\label{mag:W1-formula}
 \mathcal F_\mu(\nu)
   =\lim_{q\to\infty}\mathcal F_\mu(\nu_q)
\end{equation}
for any finite-support laws $\nu_q$ with positive rational masses such
that $W_1(\nu_q,\nu)\to0$. This limit exists and is independent of the
approximating sequence.
\end{theorem}

\begin{remark}[Conditional random spectral and field data]
\label{mag:random-data}
Let the spectrum and the field vector both be random. Conditional jointly
on these data, assume that the interaction has the Haar orbit law with
the prescribed spectrum. Let $\mu$ be a fixed compactly supported
spectral law and $\nu$ a fixed field law with finite absolute first
moment. If the spectral hypotheses of
Equation~\eqref{eq:spectral-assumptions} and convergence of the empirical
field laws to $\nu$ in $W_1$ hold in probability, then
$P_N(c^{(N)})\to\mathcal F_\mu(\nu)$ in probability. If both data
convergences hold almost surely, the pressure convergence also holds
almost surely. In either case, uniform integrability of the pressures
implies convergence of their expectations.

Indeed, condition on both data sets and apply the deterministic theorem
along convergent data sequences. The subsequence argument and localized
fourth-moment concentration in Proposition~\ref{gen:random} apply
unchanged, since the Haar Lipschitz constant is independent of the
external field. The conditional Haar hypothesis remains essential when
the fields and the interaction are dependent.
\end{remark}

\subsection{Concavity at prescribed magnetization}

We first prove the concavity needed to optimize the cavity expression.
All level derivatives in this subsection keep the interval weights
$w_i$ fixed. For a fixed bias $b$, let $Z_i$ be the tilted Gaussian
position chain, starting before the root increment at $b$, and let
$F_i$ be its unshifted backward functions. Define
\[
 \mathcal B_i(h,b)=\E\bigl[F_i'(Z_i)^2\bigr].
\]
The terminal conditional mean spin is $\tanh Z_{k-1}$, and its successive
conditional means are $F_i'(Z_i)$. Consequently these squared-mean
expectations are nonnegative and nondecreasing in $i$.

\begin{proposition}\label{mag:concavity}
Fix $s\in(-1,1)$ and positive interval weights $w_0,\ldots,w_{k-1}$.
The function $\ell^{[s]}$ is concave on
$0\le h_0\le\cdots\le h_{k-1}<\infty$. Its minimizing bias
$b_s(h)$ in \eqref{mag:constrained-functional} is unique. On strict level
data, set
\[
 B_i(h,s)=\mathcal B_i(h,b_s(h)).
\]
Then
\begin{equation}\label{mag:envelope-derivative}
 \partial_{h_i}\ell^{[s]}(h)=-\frac{w_i}{2}B_i(h,s).
\end{equation}
The associated supporting inequalities extend to boundary level data
and are unchanged by subdivision of intervals with tied covariance levels.
The function is $1/2$-Lipschitz in the $L^1$ distance between bounded
nonnegative nondecreasing paths.
\end{proposition}

\begin{proof}
Let $f_h(b)=\ell_b^{\mathrm{can}}(h)+h_{k-1}/2$ be the unshifted initial
backward value. Gaussian convolution and differentiation of the
log-exponential operators show that $f_h$ is smooth and strictly convex:
its second derivative is a positive tilted average of the terminal
second derivative, plus nonnegative conditional-variance terms. Moreover,
with $H=h_{k-1}$,
\begin{equation}\label{mag:bias-bounds}
 \log\cosh b\le f_h(b)\le\log\cosh b+H/2.
\end{equation}
For the lower bound, compare every backward operator with ordinary
expectation and use convexity of $\log\cosh$. For the upper bound, compare
with exponent one throughout and use
$\E\cosh(b+\sqrt H Z)=e^{H/2}\cosh b$.
Convex difference quotients in \eqref{mag:bias-bounds} imply that
$f_h'(b)\to\pm1$ as $b\to\pm\infty$, uniformly for $H$ in a bounded
set. Thus $f_h'(b_s(h))=s$ determines a unique bias, bounded uniformly
over paths of bounded height for each fixed interior $s$.

Gaussian differentiation, as in Proposition~\ref{fld:derivative}, gives
$\partial_{h_i}\ell_b^{\mathrm{can}}=-w_i\mathcal B_i(h,b)/2$.
The envelope rule proves \eqref{mag:envelope-derivative}. In the remainder
of the proof $s$ is held fixed, so derivatives of $B_i(h,s)$ include the
change in $b_s(h)$. We use the mean spin as spatial coordinate, so that
the varying bias is represented by the fixed initial coordinate $s$.
We will prove that $\partial_{h_l}B_i\le0$ for $l>i$ and
$\sum_l h_l\partial_{h_l}B_i\ge0$, then apply the Hessian argument of
Theorem~\ref{fld:concavity}.

Set $T=h_{k-1}$ and let $\zeta(v)=\zeta_i$ for
$h_{i-1}<v<h_i$. The continuous-time version of the unshifted backward
recursion solves, on each such interval,
\[
 F_v+\frac12F_{zz}+\frac12\zeta(v)F_z^2=0,
 \qquad F(T,z)=\log\cosh z.
\]
For the radial comparison we also need $\alpha>0$. Let $F^\alpha$ solve
the same recursion with terminal value
$\alpha^{-1}\log\cosh(\alpha z)$ and coefficient
$\alpha\zeta(v)$ in place of $\zeta(v)$. Write
\[
 u=F_z^\alpha(v,z),\qquad
 a(v,u)=F_{zz}^\alpha(v,z(v,u)),\qquad -1<u<1.
\]
Strict convexity gives the inverse $z(v,u)$. With
$L=\partial_v+\tfrac12\partial_{zz}
             +\alpha\zeta F_z^\alpha\partial_z$,
differentiation gives
$L(F_z^\alpha)=0$ and
$L(F_{zz}^\alpha)=-\alpha\zeta(F_{zz}^\alpha)^2$.
The chain rule in inverse coordinates therefore yields
\begin{equation}\label{mag:inverse-PDE}
 a_v+\frac12a^2a_{uu}=-\alpha\zeta(v)a^2,
 \qquad a(T,u)=\alpha(1-u^2).
\end{equation}
Likewise, continuation of a payoff by the tilted transition kernels
satisfies
\begin{equation}\label{mag:continuation-PDE}
 H_v+\frac12a^2H_{uu}=0.
\end{equation}
In particular, if $H_i$ has terminal condition $H_i(h_i,u)=u^2$, then
$B_i(h,s)=H_i(0,s)$ at $\alpha=1$.

We record the boundary regularity that permits comparison on the closed
interval $[-1,1]$. For fixed strict level data and $\alpha>0$, the backward
functions have, as $z\to+\infty$, expansions to every prescribed finite
order of the form
\begin{equation}\label{mag:asymptotic}
 F^\alpha(v,z)
  =z+C(v)+c_1(v)e^{-2\alpha z}+\cdots
       +c_M(v)e^{-2M\alpha z}
       +O(e^{-2(M+1)\alpha z}),\qquad c_1(v)>0.
\end{equation}
The expansions hold with any fixed finite number of differentiated
remainders, uniformly on the closed time intervals between successive
levels. To verify this, start from the expansion of the terminal
$\alpha^{-1}\log\cosh(\alpha z)$. An ordinary Gaussian backward step
integrates each exponential term. In a positive-exponent step, first
factor the linear exponential, expand the remaining exponential, integrate
term by term against the Gaussian with the resulting shifted mean, and
take the logarithm. The leading correction is multiplied by a strictly
positive Gaussian exponential moment. Remainders satisfy the same bounds:
on bounded or negative arguments an exponential bound of the indicated
order dominates their linear or exponential growth, and Gaussian
exponential moments are finite. The argument applies after differentiating
any fixed finite number of times. Coefficients are continuous in time,
and the positive leading coefficient is bounded away from zero for the
fixed data under consideration.

Tilted continuations of $(F_z^\alpha)^2$ obey the same finite expansion
property, by taking ratios of the corresponding Gaussian integrals.
Since
$1-u=2\alpha c_1(v)e^{-2\alpha z}+O(e^{-4\alpha z})$,
inverting this expansion gives bounded spatial derivatives through order
four for $a$ and the functions $H_i$ on $[-1,1]$. Reflection treats the
other endpoint. In particular,
\[
 a(v,u)=2\alpha(1-u)+O((1-u)^2)\quad(u\to1),
 \qquad
 a(v,u)=2\alpha(1+u)+O((1+u)^2)\quad(u\to-1).
\]
Thus $a$ vanishes at both endpoints, as does $(a^2)_u$.
The equations give the required one-sided time derivatives at the
finitely many level boundaries.

We next prove two monotonicities. First, every $H_i$ is convex in $u$.
Indeed $W=H_{i,uu}$ satisfies
\[
 W_v+\frac12(a^2W)_{uu}=0,
 \qquad W(h_i,u)=2.
\]
Expanded as a linear equation, its coefficients are bounded; the diffusion
and first-derivative coefficients vanish at the endpoints. Backward
comparison preserves nonnegativity. The same comparison in
\eqref{mag:continuation-PDE} shows that increasing a nonnegative
coefficient $a$ pointwise increases $H_i$, because its second spatial
derivative is nonnegative. These comparison assertions also follow
directly by reversing time, multiplying a proposed difference by a
sufficiently decaying exponential, and excluding a first negative
minimum; the bounded zeroth-order terms are absorbed by the exponential.
This argument applies interval by interval in time, including the
degenerate endpoints.

At $\alpha=1$, fix $l>i$. The covariance derivative calculation for the
backward field value, at any $v\le h_i$, reads
\[
 \partial_{h_l}F(v,z)
   =\text{a constant independent of }z
       -\frac{w_l}{2}H_l(v,F_z(v,z)).
\]
For a variation $\delta F$, differentiation at fixed inverse coordinate
gives
\begin{equation}\label{mag:inverse-variation}
 \delta a
   =\delta F_{zz}-\frac{F_{zzz}}{F_{zz}}\delta F_z
   =a^2\partial_{uu}\bigl[\delta F(v,z(v,u))\bigr].
\end{equation}
Convexity of $H_l$ therefore implies $\partial_{h_l}a\le0$ through
time $h_i$. Applying the preceding continuation comparison at the fixed
initial coordinate $s$ proves
\begin{equation}\label{mag:off-diagonal}
 \partial_{h_l}B_i(h,s)\le0\qquad(l>i).
\end{equation}

For the second monotonicity, $a$ increases with $\alpha$.
To see the sign explicitly, reverse time in \eqref{mag:inverse-PDE},
write $A_1,A_2$ for the solutions with $\alpha_1<\alpha_2$, and set
$D=A_2-A_1$. Then
\[
 D_t=\frac12A_2^2D_{uu}+cD
             +(\alpha_2-\alpha_1)\zeta A_1^2,
 \qquad
 c=(A_2+A_1)\left(\frac12A_{1,uu}+\alpha_2\zeta\right).
\]
The coefficient $c$ is bounded, the initial difference is
$(\alpha_2-\alpha_1)(1-u^2)\ge0$, and the boundary difference is zero.
Comparison proves $D\ge0$. Position and time rescaling give the identity
\[
 F_h^\alpha(v,z)
    =\alpha^{-1}F_{\alpha^2h}^{1}(\alpha^2v,\alpha z),
\]
with the level intervals rescaled on the right. The corresponding initial
biases satisfy
$b_s^\alpha(h)=\alpha^{-1}b_s^1(\alpha^2h)$.
It follows that increasing $\alpha$ increases $B_i(\alpha^2h,s)$, whence
\begin{equation}\label{mag:radial}
 \sum_l h_l\partial_{h_l}B_i(h,s)\ge0.
\end{equation}

Finally, the matrix $M_{il}=w_i\partial_{h_l}B_i$ is symmetric by
\eqref{mag:envelope-derivative}. Equations
\eqref{mag:off-diagonal} and \eqref{mag:radial} show that
$\operatorname{diag}(h)M\operatorname{diag}(h)$ has nonpositive
off-diagonal entries and nonnegative row sums. For a symmetric matrix
with these properties its quadratic form is a sum of nonnegative
weighted squares of coordinate differences and nonnegative diagonal
terms. Thus $M$ is positive semidefinite on strict level data, proving
concavity. Continuity extends the supporting inequalities to the
boundary. Tied levels introduce no new conditional mean and simply
repeat the same $B_i$. The bound $0\le B_i\le1$ in
\eqref{mag:envelope-derivative}, followed by common refinement, gives
the $L^1$ Lipschitz assertion.
\end{proof}

We shall use the weighted mean
\[
 B_i^*(h)=\sum_{j=1}^J\gamma_jB_i(h,m_j).
\]
Proposition~\ref{mag:concavity} gives the supporting derivative
$-w_iB_i^*(h)/2$ for $\ell^{[\mathbf m]}$.

\subsection{Upper bounds on magnetization slices}

We first work with a finite spectral alphabet of positive rational
proportions, along the count progressions used in the cavity construction.
We also arrange periodic field proportions $\gamma_j$, with bounded
remainder sites. The final subsection removes these restrictions.
Fix an interior rational vector $\mathbf m$. Let $\mathcal I_{j,N}$ be
the sites in field group $j$ and put $N_j=|\mathcal I_{j,N}|$. For all
sufficiently large $N$, choose integers $0\le k_{j,N}\le N_j$ such that $2k_{j,N}/N_j-1\to m_j$;
for example, take the nearest integer to $N_j(1+m_j)/2$. Define the
exact slice
\[
 \mathcal S_N(\mathbf m)=
 \left\{\sigma\in\{-1,1\}^N:
   \#\{r\in\mathcal I_{j,N}:\sigma_r=1\}=k_{j,N}
   \text{ for every }j\right\}.
\]
Retain its original cube mass $2^{-N}\ind_{\mathcal S_N(\mathbf m)}$.

The concentration, perturbation, and rotation identities of
Section~\ref{sec:rotation} apply to this prior: after normalizing its
nonzero mass, it is a deterministic probability on cube spins independent
of the Haar frame. The logarithm of its mass is kept in the pressure.
At interpolation time zero, insert an arbitrary bias $b_j$ on field
group $j$. This adds a constant on the slice. Relaxing the slice constraint
then bounds its expected enriched pressure by
\[
 \sum_j\gamma_j\ell_{b_j}^{\mathrm{can}}(h)
                -\sum_j\gamma_jb_jm_j+o(1).
\]
Minimizing in the biases gives $\ell^{[\mathbf m]}(h)+o(1)$,
uniformly on fixed-step compact field sets. Uniformity follows from
the bounded-bias conclusion of \eqref{mag:bias-bounds} and the vanishing
rounding errors in the group proportions and magnetizations.

The contact comparison of Section~\ref{sec:upper} now gives
\begin{equation}\label{mag:slice-upper}
 \limsup_{N\to\infty}\E P_{N,\mathcal S_N(\mathbf m)}(0)
      \le\inf_{p\in\cP}\{S_{\mathbf m}(p)+G_1(p)\}.
\end{equation}
Here $P_{N,\mathcal S}(0)$ denotes the pressure with the original cube
mass restricted to $\mathcal S$ and with no external field. We specify
the changes in that comparison. For a step trial $q$ of maximum less
than one, choosing bias zero in \eqref{mag:constrained-functional} and
testing $h=0$ in \eqref{mag:Sm} give
\[
 -\log2\le S_{\mathbf m}(q)\le S(q)<\infty,
\]
where finiteness follows from Lemma~\ref{fld:domain}. Replace $S(q)$
in the minimizing objective by $S_{\mathbf m}(q)$.
The time-zero bound just proved excludes that endpoint. Relaxing the
slice without any bias gives the same large-height estimate as before,
with at most an additional $\log2$ in its constant lower bound, so the
field cap is excluded as well. The field derivative at a minimum
still gives the tail inequalities between the Gibbs quantile $p$ and
$q$; the rotation identities still give the energy comparison, excluding
a minimum at positive time. These are precisely the three conclusions
used in the contact argument. The small perturbation changes the pressure
by $O(e_N)$ also for the normalized slice prior. Rounding an arbitrary
$p$ downward cannot increase $S_{\mathbf m}(p)$, and changes $G_1(p)$
continuously, proving \eqref{mag:slice-upper} for the full trial class.

It remains to justify the supremum over rational interior slices in the
full partition function. If two prescribed sets of positive-spin counts
differ by at most $\delta N$, map every configuration of the first slice
to the second by flipping at most that many spins. For a uniformly bounded
interaction matrix $J_N=U^{\mathsf T}\Lambda_NU$,
\[
 \frac1{2N}\left|\sigma^{\mathsf T}J_N\sigma
                       -\tau^{\mathsf T}J_N\tau\right|
 \le C\sqrt\delta.
\]
Indeed $\|\sigma-\tau\|\le2\sqrt{\delta N}$ and
$\|\sigma+\tau\|\le2\sqrt N$. Every image has at most the size of
a Hamming ball of radius $\delta N$ preimages. Applying maps in both
directions shows that the two restricted normalized log partitions differ
by at most $C\sqrt\delta+\mathsf h(\delta)+o_N(1)$ for small $\delta$,
where $\mathsf h(x)=-x\log x-(1-x)\log(1-x)$.
The external-field contributions differ by at most
$2\max_j|c_j|\delta$.

Choose a finite arbitrarily fine net of rational interior magnetizations,
also approximating the endpoint values $\pm1$. There are at most
$(N+1)^J$ slices, so summing their partition functions adds only
$J\log(N+1)/N$. The preceding deterministic comparison reduces their
maximum to the finite net. Concentration from
Lemma~\ref{rot:concentration} applies to every member of that net.
Equation~\eqref{mag:slice-upper}, followed by refinement of the net,
therefore proves the upper bound in \eqref{mag:finite-formula}, in
expectation and almost surely.

\subsection{A product of constrained blocks}

For the lower bound, fix $\mathbf m$ as above and choose an integer $n$
such that $n\gamma_j$, $n\gamma_j(1+m_j)/2$, and $n\rho_a$ are all
integers. Integers $n$ with these properties have unbounded multiples.
Partition the regular sites into blocks of size $n$, with the same field
group composition in each block. In every block prescribe magnetization
$m_j$ separately in group $j$. Freeze any bounded number of remainder
sites. This gives a product of nonempty finite slices; retain the
original cube mass throughout.

For this fixed $n$, add one block at a time using the coupling of
Section~\ref{sec:cavity}. The spectral coupling, fresh within-group
frames, perturbation comparison, and quadratic Gaussian transformation
are unchanged. Their base-replica moment bounds use only independence of
the fresh frames and $\|\sigma\|^2=N$, so they apply to the present prior.
For clarity, the two uses of full-model site symmetry have the following
replacements.

First, each position within a regular block is exchangeable with the
same position in every other regular block. If there are $K$ such blocks,
the expected square of a spectral projection at any fixed regular
position is at most $(N+n)/K$. For fixed $n$ this is bounded uniformly
in $N$, and the finite number of cavity coordinates consequently has the
projection tightness required for the restricted Gibbs tests. Second,
averaging a cavity spin product over its $n$ positions, against a bounded
test of the total overlap, equals the analogous average over all regular
sites. The difference from the full total overlap is bounded by twice the
fraction of remainder sites. These observations give exactly the
tightness and overlap-test identities used in
Lemma~\ref{cav:comparison} and Proposition~\ref{cav:self-consistency}.

If a restricted spin integral is empty for some frame data, assign zero
to its bounded Gibbs test. The full tilted mass outside the restriction
bounds the error of this convention. On the base side, the same moment
bounds make the probability of a nearly empty restricted mass small,
and denominator regularization proceeds as in
Lemma~\ref{cav:comparison}. For logarithms, Jensen now includes the
logarithm of the positive cavity-slice mass, a constant depending on the
fixed $n$. Thus the uniform lower bound for increments and all cap-removal
steps remain valid. No estimate uniform in $n$ is required at this point.

Let $\mathcal C_n\subset\{-1,1\}^n$ be the constrained cavity block.
Define $\ell^{(n)}(h)$ by the same backward Gaussian recursion as $\ell$,
now with $n$ independent coordinate increments and terminal payoff
\[
 \log\left(2^{-n}\sum_{\varepsilon\in\mathcal C_n}
                       e^{z\cdot\varepsilon}\right).
\]
Subtract $nh_{k-1}/2$ and divide the resulting initial value by $n$.
Let $B_i^{(n)}(h)$ be the mean over the $n$ coordinates of the expected
squared conditional spin means at level $i$. It is a nonnegative
nondecreasing sequence in $[0,1]$, and Gaussian differentiation gives
$d\ell^{(n)}=-\tfrac12\sum_iw_iB_i^{(n)}dh_i$.
For each fixed $n$, the cavity-spin tests will match these block
conditional products to the limiting overlap quantile. We then let the
block size grow. Uniform approximation of both $\ell^{(n)}$ and
$B^{(n)}$ by $\ell^{[\mathbf m]}$ and $B^*$ will allow the supporting
inequality of Proposition~\ref{mag:concavity} to control the field
supremum. Thus concavity is used for $\ell^{[\mathbf m]}$ after the
block replacement; concavity of $\ell^{(n)}$ is not required.

For a limiting base quantile $p$, take the finite monotone roundings
$\bar p$ used in Section~\ref{sec:cavity}, and set
\begin{equation}\label{mag:hp}
 h_{\bar p}(s)=\int_0^{\bar p(s)}R'(D_{\bar p}(r))\,dr.
\end{equation}
The evaluation in Proposition~\ref{cav:evaluation} gives the cavity
expression per added spin
\begin{equation}\label{mag:block-cavity}
 \ell^{(n)}(h_{\bar p})
       +\frac12\int_0^1\bar p(s)h_{\bar p}(s)\,ds+G_1(\bar p).
\end{equation}
The determinant term is unchanged. The residual covariance and the
matrix $C$ on the cavity coordinates are scalar, so their contributions
depend only on $\|\varepsilon\|^2=n$ and remain constant on the slice.

At every parameter sequence for which the averaged perturbation
estimates of Section~\ref{sec:rotation} vanish, a subsequential base
overlap limit exists. The cap argument gives a lower bound for the
liminf increment by the limit of \eqref{mag:block-cavity}. The bounded
Gibbs tests and the cavity-site average proved above give, for every
continuous $\phi$,
\[
 \lim_{\bar p\to p}\int_0^1\phi(\bar p(s))
       \bigl(B^{(n)}_{h_{\bar p}}(s)-\bar p(s)\bigr)\,ds=0.
\]
As in Proposition~\ref{cav:self-consistency}, monotonicity gives compactness
of the $B^{(n)}$ paths. A limiting path is constant on each plateau of
$p$, since equal covariance levels repeat conditional means. It is
therefore a function of $p$, and the test identity forces
\begin{equation}\label{mag:block-consistency}
 B^{(n)}_{h_{\bar p}}\longrightarrow p\quad\hbox{in }L^1.
\end{equation}
Gaussian comparison also gives the $L^1$ continuity needed for the limit
of \eqref{mag:block-cavity}. These statements take $N\to\infty$ and then
remove caps and refine the finite quantiles, with $n$ fixed.

The remaining task is to replace the block functional by the constrained
single-site functional. Its proof must be uniform in the number of
cascade levels, because the roundings of $p$ can have arbitrarily many
steps.

\begin{lemma}[Uniform replacement of a constrained block]\label{mag:replacement}
Fix an interior rational $\mathbf m$, positive rational proportions
$\gamma_j$, and a height bound $H<\infty$. Along integers $n$ for which
the block constraints are attainable,
\begin{equation}\label{mag:replacement-estimates}
 \sup_{\substack{h\in\cH\\\|h\|_\infty\le H}}
       |\ell^{(n)}(h)-\ell^{[\mathbf m]}(h)|\longrightarrow0,
 \qquad
 \sup_{\substack{h\in\cH\\\|h\|_\infty\le H}}
       \max_i|B_i^{(n)}(h)-B_i^*(h)|\longrightarrow0.
\end{equation}
The estimates are uniform in the number of steps and their positive
weights.
\end{lemma}

\begin{proof}
For a given path $h$, insert the minimizing bias $b_j=b_{m_j}(h)$ in
every coordinate of group $j$. These biases are bounded uniformly by
\eqref{mag:bias-bounds}. Consider two unnormalized $n$-coordinate
recursions: the canonical one with all $2^n$ spins, and the constrained
one with $\mathcal C_n$. Both use these same biases. Denote the backward
values just after level $i$ by $U_i$ and $V_i$, respectively, and put
$\Delta_i=U_i-V_i\ge0$. The canonical values add over coordinates, so
its tilted kernels factor over sites. Write
\[
 L_n=U_{-1}(0)-V_{-1}(0)
    =n\bigl(\ell^{[\mathbf m]}(h)-\ell^{(n)}(h)\bigr).
\]
We first prove $L_n=o(n)$ uniformly over the permitted paths, which gives
the first estimate. The same loss will bound the relative entropy of two
sampled paths; independent-site concentration under the canonical law
will then give the second estimate.
At the root, $L_n=\E\Delta_0$. At a positive-exponent step the canonical
transition kernel, denoted by $K_i^{\mathrm{can}}$, gives
\[
 \Delta_{i-1}
  =-\frac1{\zeta_i}\log
       \E_{K_i^{\mathrm{can}}}e^{-\zeta_i\Delta_i}
  \le\E_{K_i^{\mathrm{can}}}\Delta_i.
\]
Iteration proves
\begin{equation}\label{mag:root-leaf-loss}
 0\le L_n\le\E_{\mathrm{can}}\Delta_{k-1}.
\end{equation}

We bound the terminal loss without a factor depending on the number of
steps. If a one-coordinate canonical increment $g$ has Gaussian variance
$v$, its conditional tilted density is proportional to
$\exp\{\zeta F(z+g)\}$, where $F$ is $1$-Lipschitz and $0\le\zeta\le1$.
Gaussian shifting gives, for every $t\in\R$,
\begin{equation}\label{mag:increment-mgf}
 \E_{\mathrm{tilt}}e^{tg}
  =e^{t^2v/2}
    \frac{\E e^{\zeta F(z+g+tv)}}{\E e^{\zeta F(z+g)}}
  \le e^{t^2v/2+|t|v}.
\end{equation}
For the root increment the sharper ordinary Gaussian bound holds.
Successive conditioning in \eqref{mag:increment-mgf} bounds the
exponential moments of each terminal field using only total variance
$H$ and the bounded initial bias. Thus, for the vector $z$ of terminal
fields including their biases,
\begin{equation}\label{mag:field-moment}
 \E_{\mathrm{can}}\|z\|^2\le C_H n,
\end{equation}
uniformly in the depth and the step weights.

Under the averaged canonical path kernels and final spin draws, the
sites are independent and their means are $m_j$ in group $j$; this is
exactly the bias equation $f_h'(b_j)=m_j$. Map each spin vector
$\varepsilon$ deterministically to a nearest vector
$\widehat\varepsilon\in\mathcal C_n$. If $D$ is the number of flipped
coordinates, the binomial variance estimate in each group gives
$\E_{\mathrm{can}}D\le C\sqrt n$.
At fixed fields, use the pushforward of the canonical Gibbs spin law as
a trial law for the constrained log partition. The finite-space
log-sum variational identity then bounds the terminal loss by the
expected exponent loss plus the entropy loss. Averaging the former and
using \eqref{mag:field-moment},
\[
 \E\bigl|z\cdot(\varepsilon-\widehat\varepsilon)\bigr|
 \le(\E\|z\|^2)^{1/2}(4\E D)^{1/2}
 \le C_H n^{3/4}.
\]
The original spin vector is determined by its image and its flip mask.
Therefore its entropy loss is at most the entropy of that mask, and
subadditivity and concavity of binary entropy give the averaged bound
$n\mathsf h(\E D/n)\le C\sqrt n\log(n+1)$.
Together with \eqref{mag:root-leaf-loss}, this proves
\begin{equation}\label{mag:loss-rate}
 0\le L_n\le C_H\bigl(n^{3/4}+\sqrt n\log(n+1)\bigr)=o(n).
\end{equation}
This establishes the first assertion of
\eqref{mag:replacement-estimates}.

To compare the two-spin products, compare the path-mark and terminal-spin
laws under the constrained and canonical kernels. The logarithm of the
ratio of their level-$i$ densities is
$\zeta_i(\Delta_{i-1}-\Delta_i)$ for $i>0$.
The root mark law is unchanged. On allowed terminal spins, the final
log-density ratio is $\Delta_{k-1}$. Hence the chain rule for conditional
relative entropies, with weight $1/\zeta_i$ at level $i>0$ and weight one
for the final spin draw, telescopes to $L_n$:
\begin{equation}\label{mag:weighted-entropy}
 \sum_{i=1}^{k-1}\frac1{\zeta_i}
       \E_{\mathrm{con}}\KL(K_i^{\mathrm{con}}\Vert K_i^{\mathrm{can}})
 +\E_{\mathrm{con}}\KL(K_{\mathrm{spin}}^{\mathrm{con}}
                                \Vert K_{\mathrm{spin}}^{\mathrm{can}})
 =L_n.
\end{equation}
Here each conditional entropy is evaluated at the same past, with the
outer expectation taken under the constrained past law. Every term is
nonnegative. Since $\zeta_i\le1$, the ordinary relative entropy of the
full one-path laws is at most $L_n$. This bound remains valid for
arbitrarily small positive $\zeta_i$; no lower bound on the step weights
or their cumulative sums is used.

Condition now on a specified common level $i$ of two sampled cascade
leaves. The transformed cascade law and the mark representation of
Proposition~\ref{arr:rpc} show that each path has the same one-path marginal;
its kernels are shared before the split and independent afterwards.
The conditional relative entropies of the shared kernels are counted
once and those after the split, including terminal spin draws, at most
twice. Thus the joint two-path relative entropy is at most $2L_n$.
In the canonical law, sites are independent even conditional on this
branching information: the backward values add over sites and each
kernel factors over their coordinates. The average two-spin product
therefore has mean $B_i^*(h)$, and the elementary exponential-moment bound
for independent variables in $[-1,1]$ yields
\[
 \Prob_{\mathrm{can}}\left\{
   \left|\frac1n\sum_{r=1}^n\varepsilon_r^1\varepsilon_r^2-B_i^*(h)\right|
                      >\eta\right\}
     \le2e^{-n\eta^2/2}.
\]
The constants do not depend on the branching level or the depth.
For any event $A$ and probability laws $\nu,\pi$,
\[
 \KL(\nu\Vert\pi)
       \ge\nu(A)\log\frac1{\pi(A)}-\log2,
\]
as follows by mapping to the two events $A,A^c$ and expanding binary
relative entropy. Apply this with the constrained and canonical
two-path laws. Equations~\eqref{mag:loss-rate} and
\eqref{mag:weighted-entropy} imply that every fixed positive deviation
has constrained probability tending to zero uniformly in $i$ and $h$.
The average spin product is bounded by one, so its expectation converges
uniformly to $B_i^*(h)$. Its constrained expectation is $B_i^{(n)}(h)$,
which proves the second assertion.
\end{proof}

\subsection{Completing the lower bound and the approximations}

For the fixed finite spectral alphabet, $R'$ is bounded on $[0,1]$.
The fields in \eqref{mag:hp} consequently have a common height bound
$H$ independent of $p$ and its rounding. Fix a height $A<\infty$ for
a test field $v\in\cH$. By Proposition~\ref{mag:concavity} and
Lemma~\ref{mag:replacement}, after a common refinement of step intervals,
\begin{equation}\label{mag:block-tangent}
 \ell^{[\mathbf m]}(v)
 \le\ell^{(n)}(h_{\bar p})
       -\frac12\int_0^1 B^{(n)}_{h_{\bar p}}(s)
                              (v(s)-h_{\bar p}(s))\,ds
       +\epsilon_n(H,A),
 \qquad \epsilon_n(H,A)\longrightarrow0.
\end{equation}
The error is uniform over $\|v\|_\infty\le A$, the finite roundings,
and their depths. Equal covariance levels simply repeat conditional
products. Taking $\bar p\to p$ in \eqref{mag:block-tangent} and using
\eqref{mag:block-consistency} shows that the liminf cavity increment per
spin is at least
\begin{equation}\label{mag:truncated-infimum}
 a_A-\epsilon_n(H,A),\qquad
 a_A=\inf_{p\in\cP}\left\{G_1(p)+
       \sup_{\substack{v\in\cH\\\|v\|_\infty\le A}}
          \left[\ell^{[\mathbf m]}(v)+\frac12\int_0^1p(s)v(s)\,ds\right]
                         \right\}.
\end{equation}
This bound no longer depends on the limiting base array. The parameter
sets supplied by the averaged perturbation estimates have probability
tending to one. The uniform Jensen lower bound on the remaining
increments, with $n$ fixed, allows parameter averaging; telescoping along
every progression of step $n$ then gives the same lower pressure bound
for the product of constrained blocks. Perturbations change its
unnormalized expected log partition by $o(N)$, just as for the cube.

Restoring the field adds
$\sum_j\gamma_jc_jm_j+o(1)$ per spin on these configurations. They form
a subset of the original cube, so for each fixed $A$ we may first send
$N\to\infty$ at fixed admissible $n$, with the cap and rounding limits
already taken in the cavity calculation, and then send $n\to\infty$
along admissible multiples. Equation~\eqref{mag:truncated-infimum} gives
\[
 \liminf_{N\to\infty}\E P_N(c^{(N)})
        \ge\sum_j\gamma_jc_jm_j+a_A.
\]

Finally,
\begin{equation}\label{mag:height-limit}
 a_A\uparrow\inf_{p\in\cP}\{G_1(p)+S_{\mathbf m}(p)\}
                   \qquad(A\to\infty).
\end{equation}
To justify the order of optimization, choose approximate minimizers
$p_A$ along a sequence of heights tending to infinity. The set of
nondecreasing functions valued in $[0,1]$ is compact in $L^1$, so a
subsequence converges to $p$. The functional $G_1$ is continuous there.
Every fixed finite-step field $v$ is eventually allowed in the truncated
supremum, and $\int p_Av\to\int pv$. Thus the limiting truncated value
is at least $G_1(p)+\ell^{[\mathbf m]}(v)+\tfrac12\int pv$ for every
$v$. Taking the supremum proves the lower inequality in
\eqref{mag:height-limit}; the reverse inequality follows from truncation.
The values are finite above, for example by testing $p=0$, and bounded
below by testing $v=0$. This completes the lower bound for each rational
interior $\mathbf m$. Taking its supremum proves
\eqref{mag:finite-formula} under the finite rational spectral restrictions.

We remove those restrictions before passing to general fields. The count
comparison and the spectral bins in Section~\ref{sec:general} apply with
the external field present: rotations within a repeated-eigenvalue group
leave the interaction unchanged, so the same averaged projection bounds
hold. On the formula side only $G_1$ depends on the spectral law, and its
uniform difference is bounded by half the uniform difference of the
corresponding $R$ functions. Hence all the infima and the outer supremum
are stable under the spectral approximations. This proves the finite-field
formula for the compact spectral law of the theorem. Likewise, finite
field counts converging to $\gamma_j$ can be replaced, after a coordinate
permutation, by periodic rational counts with bounded remainder and
$o(N)$ changed sites. The mean absolute field change tends to zero.
Haar permutation invariance identifies the expected pressures before
and after the permutation.

For arbitrary deterministic vectors $c,d\in\R^N$, direct comparison of
the exponents gives
\begin{equation}\label{mag:field-Lipschitz}
 |P_N(c)-P_N(d)|\le\frac1N\sum_{r=1}^N|c_r-d_r|.
\end{equation}
For expected pressures we may permute coordinates first. On the line,
matching sorted vectors realizes the minimum mean absolute difference,
which is the $W_1$ distance between their equal-weight empirical laws.
Applying \eqref{mag:field-Lipschitz} to finite-field sequences and taking
their proven limits shows that $\mathcal F_\mu$ is $1$-Lipschitz on
finite-support laws with rational masses. Such laws are dense among
probability laws of finite first moment: truncate the tails, use a finite
continuity mesh, and approximate its masses by rationals. The Lipschitz
bound proves existence and independence of the limit in
\eqref{mag:W1-formula}. For a field sequence with empirical laws
$\nu_N\to\nu$ in $W_1$, fix
such an approximation $\nu_q$ and choose $N$-site field vectors with values
in $\supp\nu_q$ and proportions converging to its masses. Their empirical
laws then converge to $\nu_q$ in $W_1$. Matching the sorted vectors in
\eqref{mag:field-Lipschitz} and taking their proven finite-field limits gives
\[
 \limsup_{N\to\infty}
 \bigl|\E P_N(c^{(N)})-\mathcal F_\mu(\nu_q)\bigr|
 \le W_1(\nu,\nu_q).
\]
Letting $q\to\infty$ proves convergence of the expected pressure to
$\mathcal F_\mu(\nu)$.

Finally, the Haar Lipschitz constant of the normalized pressure is
independent of the deterministic external field. The fourth-moment
concentration estimate of Lemma~\ref{rot:concentration} and
Borel--Cantelli therefore upgrade all these expected limits to
almost-sure limits. This proves Theorem~\ref{mag:theorem}.
